\subsection{Induction, Recursion, and Counterexamples}

\begin{principle} \label{princ:mathematical-induction} \hlt{Mathematical Induction}\\
Let $P_1,P_2,\ldots$ be statements indexed by positive integers. To prove $P_n$ for every $n\in\N$, it is enough to prove:
\begin{enumerate}[label=(\roman*)]
\setlength{\itemsep}{0pt}
\item \hlt{basis}: $P_1$ is true;
\item \hlt{induction step}: for every $n\in\N$, $P_n$ implies $P_{n+1}$.
\end{enumerate}
The induction step is an implication; it does not prove $P_{n+1}$ in isolation.\\
The basis starts the chain, and the implication carries truth from one statement to the next.
\end{principle}

\begin{figure}[H]
  \centering
  \begin{tikzpicture}[
  >=Stealth,
  every node/.style={font=\small},
  value/.style={rounded corners=2pt, draw=orange!70!black, fill=orange!10, minimum width=10mm, minimum height=7mm, inner xsep=4pt, line width=0.55pt},
  arrow/.style={->, line width=0.5pt, gray!62},
  rowlabel/.style={anchor=east, gray!68, font=\small},
  rowsub/.style={anchor=east, gray!58, font=\scriptsize}
]
  \path[use as bounding box] (-2.65,-1.85) rectangle (6.25,0.52);

  \node[rowlabel] at (-0.72,0.12) {induction};
  \node[rowsub] at (-0.72,-0.18) {base + step};
  \node[value] (p0) at (0,0) {$P(0)$};
  \node[value] (p1) at (1.55,0) {$P(1)$};
  \node[value] (p2) at (3.10,0) {$P(2)$};
  \node[value] (p3) at (4.65,0) {$P(3)$};
  \node at (5.85,0) {$\cdots$};

  \draw[arrow] (p0) -- (p1);
  \draw[arrow] (p1) -- (p2);
  \draw[arrow] (p2) -- (p3);
  \draw[arrow] (p3) -- (5.55,0);

  \node[rowlabel] at (-0.72,-1.28) {recursion};
  \node[rowsub] at (-0.72,-1.58) {initial + rule};
  \node[value] (a0) at (0,-1.40) {$a_0=c$};
  \node[value] (a1) at (1.55,-1.40) {$a_1$};
  \node[value] (a2) at (3.10,-1.40) {$a_2$};
  \node[value] (a3) at (4.65,-1.40) {$a_3$};
  \node at (5.85,-1.40) {$\cdots$};

  \draw[arrow] (a0) -- (a1);
  \draw[arrow] (a1) -- (a2);
  \draw[arrow] (a2) -- (a3);
  \draw[arrow] (a3) -- (5.55,-1.40);
\end{tikzpicture}

  \caption{Induction propagates truth, while recursion propagates values along the same index chain.}
  \label{fig:peano-recursion-induction-ladder}
\end{figure}

\begin{principle} \label{princ:well-ordering-natural-numbers} \hlt{Well-Ordering of $\N$}\\
Every nonempty subset of $\N$ has a least element, justifying first unused indices, first repeated values, and increasing listings of infinite subsets of $\N$.
\end{principle}

\hlt{Placeholder}: Add strong induction, smallest-counterexample arguments, recursion, and the relation between induction and well-ordering of $\N$.
